\section{Ablation and Analysis}
\label{sec:ablations}
We examine the contribution of structured strategy knowledge, the mechanisms that construct and complete plans, and their computational cost.

\emph{Draft note:} The component scores below retain the previous 57-record study. Matched system comparisons, plan-refinement statistics, and cost measurements remain to be populated from reruns.

\subsection{From Data to Skills}
\label{sec:knowledge-ablation}
We separate three levels of agent support. \emph{Data-only} provides the input data and Bash for exploration and computation. \emph{Tool-only} adds the registered scientific tools and their existing planning and execution workflow. \emph{AgentGS} further introduces Skill organization, plan composition, and automatic completion. Data-only versus Tool-only examines the contribution of packaged scientific operations; Tool-only versus AgentGS examines the additional contribution of strategy knowledge and structured planning.

The planned comparison uses the same backbone, questions, accessible data, and model budgets. Tasks depending on data services unavailable to Bash are excluded from all three configurations. FAS provides a common outcome measure; Quality and its process components additionally compare Tool-only with AgentGS. Table~\ref{tab:system-ablation} reserves these matched results. Earlier direct-answer baselines do not substitute for the Data-only condition.
\begin{table}[t]
\caption{Planned comparison on the shared data-access scope. Numeric entries are placeholders. Quality is not applicable to Data-only (---).}
\label{tab:system-ablation}
\centering
\begin{tabularx}{\columnwidth}{@{}Xrr@{}}
\toprule
Configuration & FAS & Quality \\
\midrule
Data-only (Bash) & 0.0000 & --- \\
Tool-only & 0.0000 & 0.0000 \\
AgentGS & 0.0000 & 0.0000 \\
\bottomrule
\end{tabularx}
\end{table}


\subsection{Design Contributions}
\label{sec:design-contributions}
Within AgentGS, we consider four interventions. \emph{Flat planning} retains strategy declarations and candidate tools but asks the model to plan directly in tool instances. \emph{Without predecessor expansion} retains goal selection and dependency declarations while disabling automatic predecessor closure. \emph{Without automatic completion} disables supporting-operation insertion and bridging. \emph{Without directed bridging} removes only the search for connecting branches, retaining other completion mechanisms. These controls isolate narrower changes than removing the entire Skill layer.

Table~\ref{tab:ablation} retains the earlier results on 57 records, including 35 multistep records. Full obtains overall/multistep FAS of 0.6387/\allowbreak 0.5333. Disabling completion gives 0.4624/\allowbreak 0.3802, while disabling bridging alone gives 0.6160/\allowbreak 0.4913. The corresponding multistep decreases are 15.31 and 4.20 percentage points. These comparisons motivate examining which missing operations completion supplies and when a connecting branch changes the outcome.
\begin{table}[t]
\caption{Earlier component comparison on 57 records, including 35 multistep records. A1 is a partial control; A2 did not activate its intended intervention.}
\label{tab:ablation}
\centering
\begin{tabularx}{\columnwidth}{@{}Xrr@{}}
\toprule
Configuration & FAS & Multistep FAS \\
\midrule
Full & 0.6387 & 0.5333 \\
A1: Flat planning & 0.4748 & 0.4095 \\
A2: w/o predecessor expansion & 0.4760 & 0.4095 \\
A3: w/o automatic completion & 0.4624 & 0.3802 \\
A4: w/o directed bridging & 0.6160 & 0.4913 \\
\bottomrule
\end{tabularx}
\end{table}


Flat planning and the predecessor configuration score below Full, but their earlier runs implemented only a partial control and an inactive intervention, respectively. Their differences cannot yet establish those designs' effects. The reruns must verify effective interventions and changes in the resulting plans. Full process scores, provenance, and supplementary tool-selection and recovery comparisons are retained in Appendix~\ref{app:component-details}.

\subsection{Plan Refinement}
\label{sec:plan-refinement}
The analysis follows the plan from retrieved candidates to selected goals, expanded predecessors, and completed connections. For composition, we will report the number of candidate and retained operations, required-operation coverage, and unresolved inputs. Reference-workflow mappings and inspected traces will check whether pruning preserves necessary operations. A smaller candidate set is useful only when it retains an appropriate route to the requested outputs.

For completion, we will record how often a gap is detected, a supporting operation or branch is inserted, and grounding succeeds. Paired executions will start from the same saved initial plan with completion enabled or disabled, using matched downstream settings and budgets. We will report changes in per-task FAS, corrected and degraded outcomes, and the reasons connections are rejected. Both the full population and the subset with gaps identified before the intervention will be reported.

This comparison tests whether completion repairs an initially incomplete plan rather than benefiting from a different initial model proposal. Appendix~\ref{sec:case-study} illustrates the mechanism: a raster-producing fragment and a scalar-comparison fragment are connected through a statistics operation, whose output is explicitly bound to the comparison. The illustration explains the expected repair; its frequency and outcome effects remain to be measured.

\subsection{Quality and Cost}
\label{sec:quality-cost}
We will compare Tool-only, AgentGS, and AgentGS without completion using answer quality, model requests, token consumption, tool attempts, and elapsed time. Measurements will include failed tasks under matched budgets and execution conditions. ECR describes successful reference-call matches per attempt; actual time and token usage are needed to assess computational cost.

Comparing weaker and stronger backbones will test whether structured planning and automatic repair reduce dependence on model capability. The analysis will relate score gains to additional planning calls, inserted operations, and avoided execution failures. No efficiency advantage is inferred from the available accuracy scores alone; these measurements will determine whether improved task completion offsets the extra planning work.
